\section{Conclusion}

NoC architecture has moved to the center of AI chip design. At the foundational level, it provides the scalable packet-based communication needed to replace buses and crossbars. At the workload level, it can exploit multicast, sparsity, and mixed precision to reduce unnecessary traffic. At the system level, chiplets and heterogeneous integration transform NoC design into a package-level problem in which die boundaries, D2D standards, and compiler mapping strongly influence performance. At the trust level, the same interconnect becomes a shared attack surface that must protect sensitive model and user data. Finally, at the technology frontier, hybrid electro-photonic fabrics suggest a path beyond the bandwidth and energy limits of purely electrical links.

The common theme across these directions is that communication is no longer a background service for AI hardware. It is a first-class architectural resource whose design determines whether future AI systems can scale efficiently, securely, and economically. The strongest takeaway from our survey is therefore not that one NoC technique will dominate all future systems, but that successful AI chips will increasingly depend on communication co-design across workload, architecture, packaging, and trust.

Looking forward, several research directions appear especially promising. First, the integration of communication-aware compilation with chiplet-scale NoP design remains an open problem: compilers and runtimes need richer abstractions for die boundaries, link bandwidth, and latency so that data placement decisions can be optimized alongside routing. Second, security mechanisms must continue to evolve from static, uniform protection toward adaptive, workload-phase-aware policies that apply heavyweight protection only when and where the risk justifies the cost. Third, as photonic interconnects mature from laboratory demonstrations toward manufacturing, the architectural community needs better models for hybrid electro-photonic design-space exploration that account for thermal tuning power, insertion loss budgets, and CMOS integration cost alongside traditional NoC metrics. Finally, the trend toward open chiplet ecosystems and standardized D2D interfaces creates an opportunity for the research community to define open NoC benchmarks, traffic traces, and evaluation frameworks that can produce more consistent and comparable results across different studies.

That is the point at which NoC research stops being a narrow interconnect topic and becomes a core discipline of AI chip and system design. In a field where model sizes double every few months and deployment scenarios span from milliwatt-scale edge sensors to megawatt-scale datacenters, the communication fabric is not just infrastructure. It is the architecture.
